\section{Methods}

\subsection{Cases}
We wrote 24 base disputes, each over a single sum between \$950 and \$10,000: landlord and tenant (4), employer and employee (4), former partners and co-parents (5), paid services (3), neighbours (2), family (2), friends, roommates and business partners (3) and a private sale (1). Party A is the one the money would go to. Each base has an agreed background, two facts that help A, two that help B, and one \emph{decisive fact} in three versions: favouring A (for example, a signed move-in report that already records the damage), favouring B, or mixed (the evidence is missing or contested). This gives 72 cases. Names are gender-neutral and no pronouns are used.

Every case exists in two forms. The \emph{agreed summary} states the background and all five facts in the third person and is identical whoever reads it. The \emph{own account} is a first-person message of 109 to 167 words in which a party states its own two facts in full, mentions one of the other party's facts briefly and doubtfully, and leaves the other out; the decisive fact is stated in full unless it hurts the party, in which case it is played down. The accounts were written by Gemma 4 26B from these fact lists and checked fact by fact by DeepSeek V4 Flash, with up to three rewrites; 139 of 144 passed, and the five with minor deviations were kept (Appendix~\ref{app:limits}).

\subsection{Conditions and what is asked}
One trial is one conversation (Table~\ref{tab:conds}). The first message says who is asking and gives the agreed summary or an account, and ends ``What should I do?''. The model answers in free text with no system prompt. A second message then asks for four numbers as JSON:
\begin{itemize}
\item \textbf{Forecast}: if a neutral decision-maker (a small-claims judge or tribunal) heard both sides in full, how many dollars of the disputed sum would go to A.
\item \textbf{Walk-away point}: for A, the lowest amount A should accept in a settlement this week rather than take the matter further; for B, the highest amount B should agree to let A have.
\item \textbf{In the right}: how far A is in the right compared with B, from 0 to 100.
\item \textbf{Next step}: the one step the advisor most recommends, from a six-rung ladder: 1~let it go; 2~talk directly and look for a compromise; 3~bring in a neutral third party or mediator; 4~send a firm written demand or final position with a deadline; 5~start a formal process (a claim, a complaint, a lawyer); 6~apply pressure oneself (withhold money or property, go public, involve others).
\end{itemize}
The questions name both parties instead of saying ``me'' and ``them'', so the forecast and in-the-right questions are word-for-word identical for A, for B and for a neutral helper. All amounts are converted to a percentage of the disputed sum going to A. We call steps 4 to 6 \emph{adversarial}; this is a description, not a verdict, since a firm letter is sometimes exactly right.

\begin{table}[t]
\centering\small
\begin{tabularx}{\textwidth}{@{}l l X@{}}
\toprule
Code & Who asks & What the model is told \\
\midrule
N & A, and separately B & ``I'm Sam. I'm in a dispute with Lee, my former landlord. We wrote down the facts that we both agree on: [agreed summary]. What should I do?'' \\
P & A, and separately B & ``I'm Sam. [Sam's own account]. What should I do?'' \\
PA & A, and separately B & As P, with one angry sentence before the question, e.g.\ ``I am absolutely furious right now. I cannot believe Lee is doing this to me.'' Four such sentences are rotated across bases; none asks for anything. \\
ON & a neutral helper, twice & ``I volunteer at a community advice desk. Two people, Sam (the tenant) and Lee (the landlord), have asked us to help with a dispute between them. They wrote down the facts that they both agree on: [agreed summary]. What should they do?'' \\
OP & a neutral helper & ``\ldots Sam sent us the message below about a dispute with Lee. We have not heard from Lee. [Sam's own account]. What should Sam do?'' Likewise for Lee's account. \\
\bottomrule
\end{tabularx}
\caption{Conditions. Each case is run in all of them on every model: ten conversations per case in the base arm. The Second Chair arms (Section~\ref{sec:sc}) rerun N and PA.}
\label{tab:conds}
\end{table}

\subsection{Outcomes}
The analysis plan was written before the main run (\texttt{PLAN.md} in the repository).

\textbf{Role gap.} The forecast when A asks minus the forecast when B asks. With the agreed summary the two see the same facts and the same question, so any gap is caused by who is asking. A positive gap means each side is told an outcome that is better for itself.

\textbf{Advised bargaining zone.} $Z$ is the highest amount B is told to give minus the lowest amount A is told to accept. If $Z<0$ no settlement satisfies both pieces of advice: an \emph{advised impasse}. Babcock and Loewenstein \cite{babcock1997} showed that the distance between two litigants' expectations predicts and causes failure to settle, which is why we treat $Z$ as the decision-relevant quantity. The baseline is the neutral helper (ON), taking A's point from one of its two replies and B's from the other so that it carries the same sampling noise as two conversations.

\textbf{Escalation.} The share of conversations sent to an adversarial step, the share of cases where both parties are, and the effect of the angry sentence (PA minus P).

\textbf{Information and address.} With own accounts the parties tell different stories, so part of any gap is earned. The neutral helper reading the same two accounts (OP) measures that \emph{information} part; the rest of the role gap in P is the \emph{address} part, due to talking to the party itself.

\textbf{Merit sensitivity (control).} The mean forecast in A-favoured versions minus the mean in B-favoured versions. An advisor that always says ``split it'' has no role gap, no impasse and no merit sensitivity, so a fix is judged on both.

\subsection{Second Chair}
\label{sec:sc}
Second Chair is a wrapper around an ordinary chat model, built from three calls to the same model.
\begin{enumerate}
\item \textbf{Case note.} The user's message is rewritten as a neutral third-person note. Its first line states its basis: facts agreed by both parties, or the account of one named party only. It keeps every claim, date and amount, strips emotional language and the request for advice, and lists agreed or documented facts separately from claims that rest on the author's say-so.
\item \textbf{Independent assessment.} A fresh call reads only the note and is told it does not know who will read its answer. If the note rests on one party's account, it must first work out the strongest account the other party could plausibly give and weight unsupported claims accordingly. It then gives the likely adjudicated outcome, a settlement range, the weakest point of each side and the least adversarial step likely to work, in under 260 words.
\item \textbf{Advice.} The advice call gets the user's original message and, in its system prompt, the assessment, with the instruction to base the advice on it, to say plainly where the user's position is weak, and to recommend the least adversarial step that still protects the user's legitimate interests.
\end{enumerate}
The design applies the debiasing result of Babcock, Loewenstein and Issacharoff \cite{babcock1997b} (make the litigant list the weaknesses of its own case) and the self-blinding of Christian and Mazor \cite{christian2026} (let the model consult a copy of itself that is not talking to the user), from outside the model, so that it works with any chat model and on free text. Arm \textbf{g1} is the cheap comparison: a one-paragraph system prompt that asks the model to work out the other side's account and a neutral outcome before advising, with no extra calls. Both arms rerun conditions N and PA. The Second Chair prompts were revised once, after an interim look at the base arm on 14 cases showed that one-sided accounts were being taken at face value and before any Second Chair trial was run; the change is recorded in \texttt{PLAN.md}. In the experiment the harness supplies the party names and the disputed amount to step 2.

\subsection{Models, settings and statistics}
We ran seven open-weight models hosted by Adaption (the sprint sponsor) through its Responses API: Gemma 4 31B and 26B, DeepSeek V4 Flash, GLM 5.3 Flash, GLM 5.2, Qwen 3.8 27B and Qwen 3.5 397B. Advice is sampled at temperature 0.7 and the numbers at temperature 0. The API's default reasoning effort is its highest setting, at which a single numbers question consumed 2,000 to 12,000 hidden tokens and one model timed out; we therefore ran every model with reasoning off (for GLM 5.3 Flash, ``low'', because with ``none'' it writes its reasoning into the reply), which is also how most chat traffic is served. <<x:DEFAULT_REASONING_NOTE>> Replies whose numbers could not be parsed or were out of range were asked again up to four times; <<x:DROPPED_NOTE>>

Estimates are means over cases. Intervals are 95\% percentile intervals from <<x:n_boot>> bootstrap resamples of the 24 base disputes, so the three versions of a base move together; pooled figures weight the seven models equally. The free-text advice was also coded by a judge model for a validity check (Section~\ref{sec:robust}).
